\chapter{Transformations of Symmetry}\label{Chap:Subgroups}

When determining the structure of a material, one should not assume that when a good, or even excellent, fit to the experimental data is found that this indicates that the structural model is correct. 
There may be other models that fit better that have not yet been found. 
Even if it is known that the model offers the best fit to a set of experimental data, it may be that with more precise data, it would be clear that a different model is to be preferred. There is also the problem of "fitting noise": Sometimes, in order to improve the fit, refinement will distort a model in ways that are not chemically reasonable. If refinement improves the fit to data by distorting a phenyl group so it is no longer planar and the C-C bond lengths are unequal, that model would be far less plausible, even though it provides a better fit to the data than a chemically reasonable model. 

Perhaps a better question than "what is the correct model?" would be "what is the most accurate structure that is consistent with a set of data?" since with better data, one may be able to learn more. A key question will be to establish the best symmetry description for the structure. Crystallographers have traditionally been very concerned about missing symmetry. When a structure is refined in too low a space group, the missing symmetry will invariably degrade the quality of the refinement. Historically, the highly respected Caltech crystallographer Richard Marsh\index{Marsh, R.}, who died in 1994, published reports on crystal structures where symmetry had been missed. Invariably, the higher-symmetry structure was preferred, and often removed implausible aspects in chemistry. Dick was a pleasure to be around. I am not the only one who misses him. The process of finding a higher symmetry space group for published work is often referred to as "Marshing" in his honor. 

If the space group used for the model is too high, then the structure will be inaccurate because the symmetry will force details of the structure conform to symmetry elements that are not present, for example by requiring elements to lie on or be reflected by a mirror plane that is not actually present. When the symmetry is too high, the model will miss key details in how the structure deviates from that falsely imposed symmetry. 
When it is found that a structure is in too high a space group, there are at least a few of us that call this "Anti-Marshing". While data quality limits the ability to determine structures, my personal feeling is that using too high symmetry is actually a bigger problem than missing symmetry. As an example, reporting that a material is centrosymmetric when it is actually acentric will cause some interesting science to be missed. 

\section{Subgroup and Supergroup Relationships}\index{subgroup}\index{supergroup}

When exploring symmetry symmetry choices, it is very valuable to understand that space groups are linked together by the addition or removal of symmetry. Given a particular space group, the addition of symmetry will give rise to a \textbf{supergroup}, and the removal of a symmetry element will give rise to a \textbf{subgroup}. The relationships between space groups has been worked out via the tabulation of subgroups and supergroups. These relationships branch in that one can have a choice of several symmetry elements that can be added or removed. Likewise, these relationships form chains, since one can follow these links through multiple symmetry addition or removal steps. Note that a unit cell transformation may be involved to relate a subgroup or supergroup to the parent space group. As an example of that, consider that a mirror plane is only possible in a monoclinic or higher space group. If one removes all mirror planes from a C-centered monoclinic space group, the result will be a triclinic space group (either $P 1$ or $P \overline{1}$ to be specific) and the triclinic cell will be the primitive subcell of the C-centered monoclinic cell. 

The International Tables for Crystallography, Volume A lists the subgroups and supergroups for each space group, but not how the transformations are made, so another volume was created, International Tables Volume A1. Both volumes categorize the subgroup-supergroup relationships into categories with impressive German names, but I don't think explaining the categories is needed for this discussion. The Bilbao Crystallographic Server (BCS), at web site https://cryst.ehu.es/cryst/, has all the information in Volume A1, but with useful tools for transforming symmetry. GSAS-II performs some symmetry transformations using these tools. The BCS website will be discussed further in section \ref{ref:BCS}, below. 

Another reason why subgroup and supergroup relationships are very important has to do with what happens when a material undergoes a phase transformation. Landau theory\index{Landau, L.} (which I know very little about) describes the type of transformations that occur. One type of transformation, which Landau categorizes as \textbf{first-order}, can involve complete transformation of a material. As an example, consider the phase change between graphite, with its sheets of triangularly-bonded carbon atoms and diamond, with its 3D network of tetrahedrally bonded C atoms. Such transitions can be considered as recrystallization. However, the other type of phase transition in Landau theory, a \textbf{second-order} transition, will be a rearrangement in structure where a material's structure changes into a subgroup or supergroup. Second-order transitions are continuous and are accompanied by what Landau theory calls an order parameter, which describes how far along the transition process the transformation has gone. 

\section{Representational Analysis}\label{ref:RA}
There is an alternate approach to considering the process of phase transformations known as "representational analysis", or sometimes called "irreducible representational analysis". This is yet another subject where I am far from an expert. The history of representational analysis comes more from a spectroscopy background than from one in structural science, but in recent years there have been major advances in reconciling the nomenclature of subgroup-supergroup relationships and representational analysis. The Bilbao Crystallographic Server (section \ref{ref:BCS}) includes a number of tools for representational analysis, but Harold Stokes\index{Stokes, H.} and Dorian Hatch\index{Hatch, D.} have spent a lifetime developing tools and nomenclature for representational analysis under the collective name of ISOTROPY. Thanks to the work of Stokes, Hatch, and especially Branton Campbell\index{Campbell, B.}, as well as others, there is now a Brigham Young University website, called the ISOTROPY Software Suite, https://iso.byu.edu, which provides a number of useful tools for applying representational analysis to a structure. This is discussed further below in section \ref{ref:isotropy}.

Typical applications for representational analysis (RA) will be to describe coordinated displacement of atoms, or groupings of occupancy changes, or placing magnetic spins on atoms. Lattice strain and rotational distortions of atom groups are also possible. To provide a quick overview of RA, the symmetry of a material is classified in terms of symmetry building blocks (matrices) that cannot be further simplified, known as irreducible representations (irreps to those in the business). These irreps will be used to generate RA distortion modes which describe the relaxing the symmetry of the model by changing parameters by a formula from the value assigned to the RA mode. Representational analysis will relate two structural models, one with higher symmetry and one with lower symmetry, by a series of RA distortion modes, which each describe a set of shifts that can be applied to the model. Note that the two structural models linked by RA will be a parent-subgroup relationship, but there can be a number of subgroup links between the pair. 

Rather than vary a single atomic parameter is done in conventional crystallographic analysis, a RA mode will usually group together collective changes. For example, a RA mode might move some atoms down while moving others up. Frequently, the changes associated with an RA mode make much more chemical sense than discrete changes to crystallographic parameters, just as in spectroscopy one thinks of vibrational modes as groups of atom moving with "bends" and "stretches" rather than thinking about the individual coordinate changes. However, the number of crystallographic structural degrees of freedom will always match the total number of RA distortion modes. So, as an example if in our lower symmetry setting, we have only two atoms can be moved and those two atoms are fixed to lie on a plane, but can each be moved in two different directions, \textit{e.g.} so x and y can be changed for each atom, then there are 4 free atomic parameters, so we will also have 4 RA modes. The RA modes are formulated so that if all distortion modes are set to values of 0, the structure is in the highest symmetry setting, but as the value changes from 0, the model is reduced in symmetry. Thus, the RA modes values provide a continuous way to describe how a structure is lowered in symmetry, even though the structure will potentially move between several discrete subgroups. The Landau theory order parameter will be a vector of RA distortion modes values. 

Representational analysis provides a mechanism for exploring changes in symmetry. As well, it is useful for  categorizing potential changes in structure. However, in the end, the goal of crystallography is to describe a material's structure and for that a space group assignment is needed. 

\section{Symmetry tools in GSAS-II}

To help with determination of the optimal symmetry for a structure, GSAS-II provides a number of tools. 
Most utilize externally-created software. 
Further, there are also tools for working with distortion mode generated via representational analysis (see section \ref{ref:RA}). 
In most cases, this software is accessed via web servers, but there are also codes included with GSAS-II. 
The access to these tools is located in the section of the GUI where their use would be most convenient. 
Below is a list of the tools that GSAS-II offers for symmetry analysis, listing the 
source of the actual computations and short description of what the tool does. The presentation is organized by the section of the GUI where the tools are located.

\subsection{Histogram Unit Cells List}
The GUI that is displayed when the Unit Cells List data tree sub-entry is selected for powder diffraction histograms provides tools for indexing powder diffraction patterns and considering extinctions arising from symmetry. The upper part of the GUI in the data window is used for autoindexing and the lower portion of the window is used to superimpose the positions of reflection from a specified unit cell and space group onto the selected powder diffraction pattern. Those unit cell parameters can be loaded from indexing results, from phases that have been imported into the project (using the "Cell Index/Refine"/"Load Phase" menu command), or a cell and space group can be entered from the keyboard into the appropriate text boxes and selection pull-down menus. When a Bravais lattice has been selected and unit cell contents are entered, the reflection positions are shown as vertical dashed orange lines. The location of extinct reflections will be shown in blue when "Show Extinct" is selected. Placing the mouse on a line causes the reflection indices to be displayed as a "tooltip." The unit cell dimensions can be shifted with the arrows alongside of the text entry boxes for those dimensions and the "cell step" pull down menu dictates how much the parameters change when a arrow button is pressed. The "Try all?" button will generate a list of all the space groups consistent with the selected Bravais lattice; clicking on the entries in the list will show you the allowed reflections in that space group. 

\subsubsection{Supercell visualization}
\label{ref:supercellindexing}
Note that the unit cell parameter text entry boxes will accept math expressions. One cute trick that this allows is a simple way to test if a superlattice can index unidentified peaks. If one edits the value of \textit{a},  \textit{b}, or \textit{c} by adding "\texttt{*2}" at the end of the unit cell value, when the mouse is moved out of the box, the unit cell will be doubled. If a commensurate supercell is present, doubling \textit{a},  \textit{b}, or \textit{c} in turn should index at least some of the "extra" peaks. The original value can be restored by entering "\texttt{/2}" at the end of the new unit cell value or using the "Cell Index/Refine"/"Load Phase" menu command. 

\subsubsection{Cell Symmetry Search}
In the menu labeled "Cell Index/Refine" the "Cell Symmetry Search" command runs the NIST*LATTICE program from Vicky Karen\index{Karen, V.} and the late Alan Mighell\index{Mighell, A.}.
This program will search for unit cells of higher symmetry that are close to the input unit cell dimensions as well as possible subcells and supercells.
The generated unit cells are placed into the window where they can be compared to the histogram's powder pattern. 
The value from this is that you may find a higher symmetry cell that indexes all the powder diffraction peaks, or, a supercell may index peaks that were previously unindexed.

\subsubsection{Unit Cells List tools}
The tools that are provided for exploring symmetry in this section of the GUI use the menu labeled "Cell Index/Refine" and consist of the following menu commands: 

\begin{itemize}

    \item Run Subgroups: 
    \item Run k-SUBGROUPSMAG
    \item Transform Cell
    \end{itemize}

\subsection{Phase: General tab}
 In compute menu
\begin{itemize}
    \item Transform
    \item Std Setting
    \item Subgroups search
    \item Select magnetic/subgroup phase
    \item make subgroups project files(s)
    \item Bilbao supergroup search
    \item ISOCIF Supergroup search
\end{itemize}

\subsection{Phase: Isodistort tab}
 In compute menu
\begin{itemize}
    \item Run ISODISTORT
\end{itemize}

\subsection{Importing distortion modes from ISODISORT}

\section{The Bilbao Crystallographic Server}\label{ref:BCS}

\section{The BYU ISOTROPY Web Server}\label{ref:isotropy}


(more work needed in this chapter: \ref{needs more work})